\newpage
\pagestyle{empty} \begin{center} {\huge \textbf{감사의 글}}
\end{center}
\thispagestyle{empty}
\vspace{1cm} %{\large
% \noindent 

딥러닝 전문가라는 꿈을 이루기 위하여 석사과정을 시작하면서 많은 분들의 도움을 받아 졸업하게 되었습니다. 2년이라는 길면서도 짧은 기간 동안 연구실 생활을 하면서 저를 이끌어주시고 응원해주신 모든 분들께 감사의 마음을 전하고자 합니다.

먼저, 석사 과정 동안 지도해주신 김병규 교수님께 감사드립니다. 제가 연구하고자 하는 분야에 대하여 많은 지지를 해주시고, 방황하지 않도록 이끌어주심에 감사합니다. 졸업 논문을 준비하는 동안 자신감을 잃고 소극적으로 연구하는 저에게 야단치기 보단 묵묵히 기다려주셔서 좋은 성과를 낼 수 있었습니다. 또한 열정적으로 학문을 탐구하는 자세와 연구자가 가져야할 마음가짐을 배울 수 있었습니다. 항상 학생들의 입장에서 생각하시면서 향후 어떤 연구자가 되어야 할지 진지하게 함께 고민해주셔서 앞으로 나아가야 할 자리를 잡을 수 있었습니다.

석사 학위 논문 심사를 맡아주신 박영호 교수님과 강지우 교수님께도 감사의 말씀을 드립니다. 심사 과정 중 논문의 부족한 점을 지적해 주시고 더 나은 방향으로 이끌어 주셔서 논문을 완성할 수 있었습니다. 귀중한 시간을 내어 미래에 대한 격려와 따뜻한 조언을 해 주셔서 정말 감사드립니다.



저의 마음의 안식처이자 언제나 활기찬 연구실 생활을 만들어준 IVPL 연구실 동료분들께 정말 감사하다는 말을 전하고 싶습니다. 연구실에 어려움이 있을 때 마다 앞장서서 도와주시겠다고 하는 영운님, 항상 연구에 대하여 조언을 아끼지 않아주고 슬럼프에 빠질 때 마다 도와주는 영주언니, 

\newpage
\thispagestyle{empty}
\noindent연구실의 새벽과 주말을 함께한 메이트이자 본 받을 점이 많은 혜진언니, 묵묵히 연구하면서 재치있게 분위기를 풀어주는 정인이, 연구실 생활동안 함께 과제를 진행하며 성장할 수 있도록 도와준 예지언니, 연구실의 막내이지만 언니들 보다 더 언니같아 배울 점이 많던 한림이까지 모두들 감사합니다. 비록 제가 석사 생활을 시작했을 때 바로 졸업을 하여 함께 연구실 생활을 한 기간은 짧지만, 졸업 이후에도 간간히 연락을 하면서 항상 응원해준 영진이와 운하에게도 감사하다는 말을 전합니다. 

언제나 아낌없는 지원과 응원을 해준 가족과 친구들에게도 감사를 전합니다. 석사 과정을 하면서 항상 멋있다고 격려해주시는 아버지, 어려움이 있을 때 실질적 해결책을 주시려는 어머니, 친구같은 둘째 동생 예린이, 그리고 가장 사랑하는 막내 리나 언제나 응원하고 위로해 주어서 감사합니다. 나의 가장 친한 친구이자 힘들때 마다 격려해준 채린, 존재만으로 든든한 살리고 멤버들 주현, 소연 그리고 수진이에게 감사드립니다. 제주도 메이트이자 본격적인 프로그래밍 공부를 시작하는데 도움을 많이 준 소희, 유진, 민지 그리고 수빈에게도 감사드립니다. 그리고 졸업프로젝트할때 카메라를 빌려주고 대학원 생활에 대하여 소소한 팁들을 남겨준 강타님께도 감사합니다.

많은 사람들의 지지가 있었기에 학위과정을 포기하지 않고 해낼 수 있었다고 생각합니다. 석사 학위 과정 동안 배운 지식과 터득한 지혜로 사회에 도움이 되는 사람이 되도록 노력하겠습니다. 감사합니다.